\itemhead{习题4.2 · 第1题（教材p.113）}{XJ-B1-C04-HW-P113-Y0886}
\label{XJ-B1-C04-HW-P113-Y0886}
{\footnotesize 湘教版教材p.113 / PDF第120页 / 习题4.2 / 第1题}
\field{原题}{设 $f(x)=2^x$，计算 $f(0)-f(-1)$，$f(2)-f(1)$，$f(4)-f(3)$，$f(6)-f(5)$。比较这些差的大小，你有什么发现？}
\field{参考答案}{四个差依次为 $1/2,2,8,32$，逐渐增大；一般 $f(n+1)-f(n)=2^n$，相邻单位区间的增长量随 $n$ 增大而增大。}
{\small \field{关联知识}{E01 指数函数的概念、图象与性质（教材帮p.150、151 / PDF159、160）；P02 有理数、实数指数幂与运算（教材帮p.144、145 / PDF153、154）}
\field{方法与易错点}{先查底数条件，利用已知点定底数，再据单调性及图象求值或范围。 易错：只要求底数非负；漏排1；将 $a^x+b$ 当标准指数函数且仍用原值域。}
\field{答案依据}{答案册相关题组：PDF第32页（册内页31）；按当前题库原题逐项重算，不按旧版题号机械对套。}}
